\chapter{高等代数：用正交特征基证明全等}
\label{ch:linear}

对称性是谱定理的关键。若 $A^T=A$，则 $\langle Ax,y\rangle=\langle x,Ay\rangle$。若 $Av=\lambda v,\ Aw=\mu w$ 且 $\lambda\ne\mu$，这个等式给出 $(\lambda-\mu)\langle v,w\rangle=0$，即两个特征方向垂直。但这还没有保证特征向量够用，下面给出完整论证。

\begin{theorem}[实对称矩阵的谱定理]
每个实对称矩阵 $A$ 都有一组标准正交的特征基。把它们作为列组成 $Q$，便有 $Q^TAQ=D$，且 $Q^TQ=I$。
\end{theorem}
\begin{proof}
单位球面紧致，连续函数 $x^TAx$ 在其上取得最大值。设单位向量 $v$ 是一个最大点。任取单位向量 $w\perp v$，函数
\[
t\longmapsto(\cos t\,v+\sin t\,w)^TA(\cos t\,v+\sin t\,w)
\]
在 $t=0$ 的导数是 $2w^TAv$，必须为零。因此 $Av$ 垂直于 $v^\perp$，即 $Av=\lambda v$。对任意 $z\perp v$，有 $v^TAz=(Av)^Tz=\lambda v^Tz=0$，故 $A$ 把 $v^\perp$ 映回自身，且限制在其中仍是对称的。在这个维数少一的空间重复论证；一维时结论直接成立。这就得到所需的正交特征基。
\end{proof}

\section{从系数算出主轴和距离}

\begin{example}{从原方程开始算出旋转矩阵}
不借用导读的答案，计算 $5x^2+4xy+2y^2$ 的正交标准形，并验证两对对应点之间的距离。
\end{example}
\begin{solution}
交叉项的系数 $4$ 分给矩阵的两个非对角位置，各为 $2$。于是
\[
\det(A_0-\lambda I)=(5-\lambda)(2-\lambda)-4
=(\lambda-6)(\lambda-1).
\]
当 $\lambda=6$ 时，$-x+2y=0$，可取方向 $(2,1)$；当 $\lambda=1$ 时，$2x+y=0$，可取 $(-1,2)$。两者的长度均为 $\sqrt5$，除以长度得到 $Q_0$ 的两列。由于 $A_0Q_0=Q_0D_0$ 且 $Q_0^TQ_0=I$，得到标准形 $6u^2+v^2$。

例如在标准椭圆取 $P=(1/\sqrt6,0)$ 和 $R=(0,1)$，距离平方是 $1/6+1=7/6$。它们的像是 $P'=(2,1)/\sqrt{30}$、$R'=(-1,2)/\sqrt5$；把坐标差平方相加，交叉项相消，仍为 $7/6$。这个计算核验了一对点；两椭圆全等的证明需要对任意 $p,r$ 使用 $\|Q_0(p-r)\|=\|p-r\|$，再用 $Q_0^{-1}$ 证明一一对应。
\end{solution}

\begin{example}{负特征值不妨碍等值图形全等}
把 $4xy=1$ 化成无交叉项的方程，并说明换元后得到的是什么图形。
\end{example}
\begin{solution}
系数矩阵 $B=\begin{pmatrix}0&2\\2&0\end{pmatrix}$。向量 $(1,1)^T$ 与 $(-1,1)^T$ 的特征值分别为 $2,-2$。令
\[
Q=\frac1{\sqrt2}\begin{pmatrix}1&-1\\1&1\end{pmatrix},
\qquad
x=\frac{u-v}{\sqrt2},\quad y=\frac{u+v}{\sqrt2}.
\]
代入得 $4xy=2u^2-2v^2$。因此原方程与 $2u^2-2v^2=1$ 的双曲线全等。后者的顶点是 $(\pm1/\sqrt2,0)$，送回原坐标为 $\pm(1/2,1/2)$；渐近线 $u=\pm v$ 对应原来的两条坐标轴。正负号决定图形的类型；能否全等对应由 $Q$ 是否正交决定。
\end{solution}

\section{重特征值与投影}

\begin{example}{在二维特征空间中找两条正交轴}
设 $w=(1,1,1)^T$，$B=I+ww^T$。求正交标准形，并找出 $x^TBx=1$ 的三个半轴长度。
\end{example}
\begin{solution}
由于 $w^Tw=3$，有 $Bw=w+w(w^Tw)=4w$。若 $w^Tx=0$，则 $Bx=x$，所以垂直于 $w$ 的平面是特征值 $1$ 的特征空间。取
\[
v_1=\frac{(1,1,1)^T}{\sqrt3},\quad
v_2=\frac{(1,-1,0)^T}{\sqrt2},\quad
v_3=\frac{(1,1,-2)^T}{\sqrt6}.
\]
逐项内积可验证正交且都是单位向量。取这些列得到 $Q$，于是 $Q^TBQ=\operatorname{diag}(4,1,1)$。对应椭球方程 $4u_1^2+u_2^2+u_3^2=1$，三个半轴长为 $1/2,1,1$。后两根轴可在同一平面内一起转动；不变的是该平面与半径，而不是某两根指定的轴。
\end{solution}

\begin{example}{不选特征向量的正负号，也能投影}
只用 $A_0$ 的矩阵与特征值 $6,1$，求投影到特征值 $6$ 的正交投影矩阵。将 $e=(1,0)^T$ 分成两个正交分量，计算 $A_0^2e$。
\end{example}
\begin{solution}
多项式 $p(t)=(t-1)/5$ 在 $6$ 处取 $1$、在 $1$ 处取 $0$。因此 $p(A_0)$ 在两个特征方向上分别乘 $1,0$，故
\[
P_6=\frac{A_0-I}{5}=\frac15\begin{pmatrix}4&2\\2&1\end{pmatrix},
\qquad
P_1=I-P_6=\frac15\begin{pmatrix}1&-2\\-2&4\end{pmatrix}.
\]
$P_6e=(4,2)^T/5$，$P_1e=(1,-2)^T/5$，两者内积为零，和为 $e$。作用两次等于把这两个分量分别乘 $36,1$，得到
\[
A_0^2e=36P_6e+P_1e=(29,14)^T.
\]
这与直接求 $A_0(5,2)^T$ 一致。投影矩阵不随单位特征向量的正负号改变；后面将用复积分和范畴论再得到同一投影。
\end{solution}

\section{什么时候可以用谱判定全等}

\begin{example}{反过来：全等的正定单位椭圆一定同谱吗}
设 $A,B$ 都是实对称正定的二阶矩阵。证明 $x^TAx=1$ 和 $y^TBy=1$ 全等，当且仅当 $A,B$ 的特征值（计重数）相同。
\end{example}
\begin{solution}
若两矩阵同谱，把各自的单位特征基按相同顺序排列，写成 $A=Q_ADQ_A^T$、$B=Q_BDQ_B^T$。则 $U=Q_BQ_A^T$ 正交，且 $U^TBU=A$，所以 $y=Ux$ 给出保距对应。

反过来，任一平面刚体等距变换写成 $y=Ux+b$，其中 $U$ 正交。椭圆的中心唯一：若一个非空有界点集有两个中心，则两次中心对称的复合是非零平移；反复平移会产生无界点列，矛盾。两椭圆中心都是原点，故 $b=0$。任取 $v\ne0$，令 $x=v/\sqrt{v^TAv}$，则 $x$ 在第一条椭圆上，故 $(Ux)^TB(Ux)=1$。乘回分母得到 $v^TAv=v^TU^TBUv$。这对所有 $v$ 成立；分别令 $v=e_i$、$e_i+e_j$，可读出两边对称矩阵的各项系数相等，所以 $A=U^TBU$。正交相似保特征值。
\end{solution}

\begin{example}{同谱的非对称矩阵能否照搬这个结论}
比较 $C=I_2$ 和 $J=\begin{pmatrix}1&1\\0&1\end{pmatrix}$。它们是否同谱、相似？对应的实二次型单位等值集是否全等？
\end{example}
\begin{solution}
两者的特征多项式都是 $(\lambda-1)^2$，但 $J-I$ 的核只有一维，$C-I$ 的核是二维。相似保持核的维数，因此它们不相似，当然也不正交相似。另一个直接理由是任意可逆 $S$ 都有 $S^{-1}IS=I$。

对实向量 $x$，$x^TJx=x^T\bigl((J+J^T)/2\bigr)x$，因为反对称部分在两侧乘同一向量时为零。$J$ 的对称部分为 $\begin{pmatrix}1&1/2\\1/2&1\end{pmatrix}$，特征值是 $3/2,1/2$。它的单位等值椭圆半轴长为 $\sqrt{2/3}$ 和 $\sqrt2$；$C$ 的单位等值集是半径 $1$ 的圆。由上一题可知不全等。对称性不能省；一般实矩阵的二次型由对称部分决定。
\end{solution}
